\documentclass[a4paper,10pt]{report}\input{../0.Préambule/Préambule_cours_prof}\begin{document}

\newpage
\pagestyle{empty}
Troisième Générale \hfill Session Janvier 2022

\vspace{1cm}
\begin{minipage}{8cm}
Nom : \dotfill
\end{minipage}
\hspace{5mm}
\begin{minipage}{8cm}
Prénom : \dotfill
\end{minipage}

\vspace{1cm}
\noindent\hrulefill
\begin{center}
\section*{\hspace{3.5cm} Devoir commun de Mathématiques \hspace{3.5cm}}
\bigbreak
\subsection*{Géométrie - Programmation - Statistiques - Calcul algébrique - Arithmétique}
\noindent\hrulefill
\end{center}

\vspace{0.5cm}

\begin{center}
\subsection*{\hspace{4.5cm} Durée de l'épreuve : 2 heures\hspace{4.5cm} }
\end{center}

\vspace{0.5cm}


\noindent \textbf{Le sujet est composé de $5$ pages.}

\medskip
\begin{enumerate}
\item[•] Toute réponse doit être parfaitement justifiée (sauf indication contraire dans la consigne).
\item[•] La qualité, la clarté de la rédaction et la maîtrise de la langue seront prises en compte dans l'appréciation de la copie.
\item[•] Les exercices seront traités dans l'ordre de votre choix.
\item[•] L'emploi de la calculatrice est autorisé.
\item[•] \textbf{L'énoncé est à rendre avec la copie.}( Exercices $4$ et $7$ à compléter sur le sujet).
\end{enumerate}

\vspace{1cm}

\textbf{Barème des exercices :}

\begin{center}
\begin{minipage}{10cm}
\renewcommand{\arraystretch}{1.5}
\begin{tabularx}{\linewidth}{|*{4}{>{\centering \arraybackslash}X|}}
\hline
\rowcolor{gray!20}Exercice & \multicolumn{2}{|c|}{Nombre de points} & Compétences \\\hline
$1$ & & $8$ & Calculer \\\hline
$2$ & & $12$ & Calculer \\\hline
$3$ & & $10$ & Chercher\\\hline
$4$ & & $16$ & \\\hline
$5$ & & $14$ &  \\\hline
$6$ & & $28$ & Raisonner\\\hline
$7$ & & $12$ & Modéliser\\\hline
Total & & $100$ & \\\hline
\end{tabularx}
\end{minipage}
\end{center}

\vspace{1cm}

\begin{center}
Collège Charles DULLIN
\end{center}

\newpage
\paragraph{Exercice 1}: \hfill \textbf{8 points}
\medskip

Pour chacun des calculs ci dessous, on écrira toutes les étapes de calcul. Il ne sera attribué aucun points si la réponse n'est pas détaillée

\begin{enumerate}
\item Calculer $A$ et donner le résultat sous forme de fraction irréductible. $$A = \dfrac{5}{3}- \dfrac{4}{3} \times \dfrac{7}{5}$$

\item Calculer $B$ et donner le résultat en écriture scientifique. $$B=\dfrac{210 \times 10^{-6} \times 5 \times 10^{5}}{35 \times 10^{4}}$$
\end{enumerate}

\paragraph{Exercice 2}: \hfill \textbf{12 points}
\medskip

\begin{enumerate}
\item Monter que  la décomposition en produit de facteurs premier de $288$ et $224$ est : 

$$288 = 2^5 \times 3^2 \quad \quad \quad \quad 224 = 2^5 \times 7$$
\item Calculer le plus grand diviseur commun à  $288$ et à $224$.
\item Écrire la fraction $\dfrac{288}{224}$ sous forme irréductible.
\item Un photographe doit réaliser une exposition en présentant ses œuvres sur des panneaux contenant chacun le même nombre de photos de paysages et le même nombre de portraits. Il dispose de $224$ photos de paysages et de $288$ portraits.
\begin{enumerate}
\item Combien peut-il réaliser au maximum de panneaux en utilisant toutes les photos ?
\item Combien chaque panneau contient-il de paysages et de portraits ?
\end{enumerate}
\end{enumerate}

\paragraph{Exercice 3}: \hfill \textbf{10 points}
\medskip

Maxime et Claire sont célibataires, n'ont pas d'enfant et paient leur impôt séparément.

\noindent Auraient-ils payé moins d'impôt s'ils étaient mariés ( en payant ensemble un seul impôt) ?

\noindent Expliquer clairement la démarche, en vous appuyant sur les documents joints et en détaillant les calculs.

\noindent \begin{minipage}{6cm}
\begin{doc}[Revenus imposables :]

\begin{enumerate}
\item[•] Maxime : $\np{25905}$\euro
\item[•] Claire : $\np{27750}$\euro
\end{enumerate}
\end{doc}
\end{minipage}
\hspace{0.5cm}
\begin{minipage}{10.5cm}
\begin{doc}[Nombre de parts $N$ suivant la situation familiale :]

\begin{tabularx}{\linewidth}{|*{2}{>{\centering \arraybackslash}X|}}
\cline{2-2}
\multicolumn{1}{c|}{}& \cellcolor{doc!20}{Nombres de parts} \\\hline
\cellcolor{doc!20}{Célibataire sans enfant} & $1$ \\\hline
\cellcolor{doc!20}{Couple sans enfant} & $2$ \\\hline
\end{tabularx}
\end{doc}
\end{minipage}

\noindent \begin{minipage}{8cm}
\begin{doc}[Barème de l'impôt sur le revenu :]

$R$ désigne le montant du revenu imposable (en euros) et $N$ le nombre de parts.

\begin{tabularx}{\linewidth}{|*{2}{>{\centering \arraybackslash}X|}}
\hline
\cellcolor{doc!20}{Tranche du revenu net imposable (en euros)} & \cellcolor{doc!20}{Formule de calcul de l'impôt} \\\hline
jusqu'à $\np{9700}$& $-$ \\\hline
de $\np{9701}$ à $\np{26791}$& $0,14 \times R - \np{1358} \times N$ \\\hline
de $\np{26792}$ à $\np{71826}$& $0,3 \times R - \np{5644,56} \times N$ \\\hline
de $\np{71827}$ à $\np{152108}$& $0,41 \times R - \np{13545,42} \times N$ \\\hline
plus de $\np{152108}$& $0,45 \times R - \np{19629,74} \times N$ \\\hline
\end{tabularx}

\end{doc}
\end{minipage}
\hspace{0.5cm}
\begin{minipage}{8.5cm}
\begin{doc}[Aide au calcul]

\begin{enumerate}
\item[•]Le revenu imposable est l'ensemble des revenus, bénéfices et gains perçus pendant l'année.
\item[•]Le nombres de parts est un nombre qui intervient dans le calcul de l'impôt et que dépend de la situation familiale (mariés, célibataire, avec ou sans enfant, ...)
\item[•]Pour un couple, le revenu imposable est la somme des revenus de chacun des partenaires.
\end{enumerate}
\end{doc}
\end{minipage}

\newpage
\paragraph{Exercice 4}: \hfill \textbf{16 points}
\medskip

Dans cet exercice, des traits sur le graphique pourront servir de justification au différentes questions.

\medskip
Une famille a effectué une randonnée en montagne. Le graphique ci-dessous donne la distance parcourue en km en fonction du temps en heures.

\begin{center}
\begin{tikzpicture}[line cap=round,line join=round,>=triangle 45,x=1.5cm,y=0.35cm]
\begin{axis}[x=1.5cm,y=0.35cm,axis lines=middle,ymajorgrids=true,xmajorgrids=true,
xmin=0.0,xmax=7.1,ymin=0.0,ymax=22.0,xtick={0.0,1.0,...,7.0},ytick={0.0,1.0,...,22.0},]
\clip(-5,-5) rectangle (7.1,25.);
\draw [line width=1.2pt] (0.,0.)-- (1.,4.);
\draw [line width=1.2pt] (1.,4.)-- (2.,7.);
\draw [line width=1.2pt] (2.,7.)-- (3.,8.);
\draw [line width=1.2pt] (3.,8.)-- (4.,15.);
\draw [line width=1.2pt] (4.,15.)-- (5.,15.);
\draw [line width=1.2pt] (5.,15.)-- (6.,18.);
\draw [line width=1.2pt] (6.,18.)-- (7.,20.);
\end{axis}
\node[below=5mm] at (3.5,0) {Temps (en heures)};
\node[left=12mm,rotate=90] at (0,20) {Distance parcourue (en km)};
\end{tikzpicture}
\end{center}

\begin{enumerate}
\item Ce graphique traduit-il une situation de proportionnalité? Justifier la réponse.
\item On utilisera le graphique pour répondre aux questions suivantes.
\begin{enumerate}
\item Quelle est la durée totale de cette randonnée ?
\item Quelle distance cette famille a-t-elle parcourue au total ?
\item Quelle est la distance parcourue au bout de $6$h de marche ?
\item Au bout de combien de temps ont-ils parcouru les $8$ premiers km?
\item Que s'est-il passé entre la $4^e$ et la $5^e$ heure de randonnée ?
\end{enumerate}
\item Un randonneur expérimenté marche à une vitesse moyenne de $4$km/h sur toute la randonnée. Cette famille est-elle expérimentée ? Justifier la réponse.
\end{enumerate}

\paragraph{Exercice 5}: \hfill \textbf{14 points}
\medskip

Voici les résultats du brevet blanc de deux classes de $3^e$ du collège POUDLARD .

\medskip
\noindent \textbf{GRYFFONDOR :}

$8$\quad ; \quad $7$\quad ; \quad $12$\quad ; \quad $15$\quad ; \quad $15$\quad ; \quad $12$\quad ; \quad $18$\quad ; \quad $18$\quad ; \quad $11$\quad ; \quad $7$\quad ; \quad $8$\quad ; \quad $11$\quad ; \quad $7$\quad ; \quad $13$\quad ; \quad $10$\quad ; \quad $10$\quad ; \quad $6$\quad ; \quad $11$

\medskip
\noindent \textbf{SERPENTARD} : Le professeur Rogue à déjà rentré les notes de Serpentard dans un tableur :

\begin{center}
\begin{tableur}{11}
\hline
& Notes & $7$ & $8$ & $9$ & $12$ & $13$ & $16$ & $18$ & $19$ & Total &\\\hline
& Effectifs & $4$ & $3$ & $2$ & $1$ & $4$ & $1$ & $2$ & $1$ &  &\\\hline
\end{tableur}
\end{center}

\begin{enumerate}
\item Quelle formule doit-on entrer en $J2$ ? Calculer ce total.
\item Calculer la moyenne de chacune de ces classes (arrondir au dixième).
\item Calculer pour chaque classe, le pourcentage d'élèves qui ont la moyenne (arrondir à l'unité).

Que constatez-vous ?
\end{enumerate}

\newpage
\paragraph{Exercice 6}: \hfill \textbf{28 points}	
\medskip

Dans l'exercice suivant, les figures ne sont pas à l'échelle.

\medskip
\begin{minipage}{5cm}
\begin{center}
\begin{tikzpicture}[line cap=round,line join=round,>=triangle 45,x=0.5cm,y=0.5cm]
\clip(0.9,-4.1) rectangle (8.1,10.6);
\fill[line width=1.2pt,fill=black,fill opacity=0.20000000298023224] (1.,7.) -- (8.,7.) -- (8.,6.5) -- (1.,6.5) -- cycle;
\fill[line width=1.2pt,fill=black,fill opacity=0.20000000298023224] (1.,3.5) -- (8.,3.5) -- (8.,3.) -- (1.,3.) -- cycle;
\fill[line width=1.2pt,fill=black,fill opacity=0.20000000298023224] (1.,0.) -- (8.,0.) -- (8.,-0.5) -- (1.,-0.5) -- cycle;
\fill[line width=1.2pt,fill=black,fill opacity=0.20000000298023224] (1.,10.5) -- (8.,10.5) -- (8.,10.) -- (1.,10.) -- cycle;
\fill[line width=1.2pt,fill=black,fill opacity=0.20000000298023224] (1.,-3.5) -- (8.,-3.5) -- (8.,-4.) -- (1.,-4.) -- cycle;
\draw [line width=1.2pt] (5.,10.)-- (1.,7.);
\draw [line width=1.2pt] (1.,10.)-- (7.,7.);
\draw [line width=1.2pt] (1.,7.)-- (8.,7.);
\draw [line width=1.2pt] (8.,7.)-- (8.,6.5);
\draw [line width=1.2pt] (8.,6.5)-- (1.,6.5);
\draw [line width=1.2pt] (1.,6.5)-- (1.,7.);
\draw [line width=1.2pt] (1.,3.5)-- (8.,3.5);
\draw [line width=1.2pt] (8.,3.5)-- (8.,3.);
\draw [line width=1.2pt] (8.,3.)-- (1.,3.);
\draw [line width=1.2pt] (1.,3.)-- (1.,3.5);
\draw [line width=1.2pt] (1.,0.)-- (8.,0.);
\draw [line width=1.2pt] (8.,0.)-- (8.,-0.5);
\draw [line width=1.2pt] (8.,-0.5)-- (1.,-0.5);
\draw [line width=1.2pt] (1.,-0.5)-- (1.,0.);
\draw [line width=1.2pt] (1.,10.5)-- (8.,10.5);
\draw [line width=1.2pt] (8.,10.5)-- (8.,10.);
\draw [line width=1.2pt] (8.,10.)-- (1.,10.);
\draw [line width=1.2pt] (1.,10.)-- (1.,10.5);
\draw [line width=1.2pt] (1.,-3.5)-- (8.,-3.5);
\draw [line width=1.2pt] (8.,-3.5)-- (8.,-4.);
\draw [line width=1.2pt] (8.,-4.)-- (1.,-4.);
\draw [line width=1.2pt] (1.,-4.)-- (1.,-3.5);
\draw [line width=1.2pt] (1.,10.)-- (1.,-3.5);
\draw [line width=1.2pt] (1.,6.5)-- (7.,3.5);
\draw [line width=1.2pt] (5.,6.5)-- (1.,3.5);
\draw [line width=1.2pt] (5.,3.)-- (1.,0.);
\draw [line width=1.2pt] (1.,3.)-- (7.,0.);
\draw [line width=1.2pt] (5.,-0.5)-- (1.,-3.5);
\draw [line width=1.2pt] (1.,-0.5)-- (7.,-3.5);
\end{tikzpicture}

Plateau en bois\\
d'épaisseur $2$cm

\medskip
 Figure 1
\end{center}
\end{minipage}
\hspace{1cm}
\begin{minipage}{10cm}
Un décorateur a dessiné une vue de coté d'un meuble de rangement composé d'une structure métallique et de plateaux en bois d'épaisseur $2$cm, illustré par la figure 1.

\medskip
\noindent Les étages de la structure métallique de ce meuble de rangement sont tous identiques et la figure $2$ représente l'un d'eux.

\begin{center}
\begin{tikzpicture}[line cap=round,line join=round,>=triangle 45,x=1.0cm,y=1.0cm]
\clip(0.5,1.5) rectangle (6.5,5.5);
\draw [line width=1.2pt] (1.,5.)-- (1.,2.);
\draw [line width=1.2pt] (1.,5.)-- (4.,5.);
\draw [line width=1.2pt] (4.,5.)-- (1.,2.);
\draw [line width=1.2pt] (1.,5.)-- (6.,2.);
\draw [line width=1.2pt] (6.,2.)-- (1.,2.);
\begin{scriptsize}
\draw[color=black] (0.7664693243515834,5.278328752847963) node {\large{$A$}};
\draw[color=black] (0.7386092123247671,1.795814749495913) node {\large{$C$}};
\draw[color=black] (4.304703551757265,5.292258808861371) node {\large{$B$}};
\draw[color=black] (6.254911393634412,1.7818846934825046) node {\large{$D$}};
\draw[color=black] (2.883837838389629,3.551001807185346) node {\large{$O$}};
\end{scriptsize}
\end{tikzpicture}

Figure 2
\end{center}
\end{minipage}

\medskip
On donne :
\begin{enumerate}
\item[•] $OC = 48$cm; $OD = 64$cm ; $OB = 27$cm ; $OA = 36$cm et $CD = 80$cm;
\item[•] les droites $(AC)$ et $(CD)$ sont perpendiculaires.
\end{enumerate}

\medskip
\begin{enumerate}
\item Démontrer que les droites $(AB)$ et $(CD)$ sont parallèles.
\item Montrer par le calcul que $AB = 45$cm.
\item Calculer la hauteur totale du meuble de rangement.
\end{enumerate}

\paragraph{Exercice 7}: \hfill \textbf{12 points}
\medskip

Une partie de l'exercice est à compléter sur le sujet.
\medskip

Un supermarché propose une animation pour les clients :

\noindent Quand un client entre dans le magasin, un petit jeu sur tablette est présenté sous forme du programme scratch (avec comme lutin un chat) donné en \textbf{ANNEXE}

\begin{enumerate}
\item Compléter le tableau d'étape suivant avec $2$ comme nombre choisi au départ ?

\renewcommand{\arraystretch}{2}
\arrayrulecolor{black}
\begin{tabularx}{\linewidth}{|p{2cm}|*{8}{>{\centering \arraybackslash}X|}}
\hline
Ligne & $1$ & $2$ & $3$ & $4$ & $5$ & $6$ & $7$ & $8$ \\\hline
Réponse & \cellcolor{gray!50}{X} & \cellcolor{gray!50}{X} &  &  &  &  &  &  \\\hline
Nombre & \cellcolor{gray!50}{X} & \cellcolor{gray!50}{X} & \cellcolor{gray!50}{X} &  &  &  &  &  \\\hline
\end{tabularx}

\medskip
Que dira le chat à l'issue du programme ?
 
\item Que dira le chat si la valeur choisie au départ est $7$ ?
Justifier la réponse.
\item En notant $x$ le nombre de départ, montrer que le programme calcule simplement le triple du nombre choisi au départ.
\end{enumerate}

\newpage
\begin{minipage}{8cm}
\paragraph{ANNEXE}: à coller sur votre copie
\medskip

\begin{scratch}[num blocks]
\blockinit{quand \greenflag est cliqué}
\blockmove{demander \txtbox{bonjour, choisissez un nombre} et attendre}
\blockvariable{mettre {\selectmenu{nombre} à \boolsensing{réponse}}}
\blockvariable{mettre {\selectmenu{nombre} à \ovaloperator{\ovalvariable{nombre} * \ovalnum{1,5}}}}
\blockvariable{ajouter à \selectmenu{nombre} \ovalnum{3}}
\blockvariable{mettre \selectmenu{nombre} à \ovaloperator{\ovalvariable{nombre} * \ovalnum{2}}}
\blockvariable{mettre \selectmenu{nombre} à \ovaloperator{\ovalvariable{nombre} - \ovalnum{6}}}
\blockifelse{si \booloperator{\ovalvariable{nombre} < \ovalnum{20}} alors}
{\blocklook{dire \ovalnum{Bravo, vous avez gagné} pendant \ovalnum{2} seconde}}
{\blocklook{dire \ovalnum{Désolé, vous avez perdu} pendant \ovalnum{2} seconde}}
\end{scratch}
\end{minipage}
\hspace{5mm}
\begin{minipage}{8cm}
\paragraph{ANNEXE}: à coller sur votre copie
\medskip

\begin{scratch}[num blocks]
\blockinit{quand \greenflag est cliqué}
\blockmove{demander \txtbox{bonjour, choisissez un nombre} et attendre}
\blockvariable{mettre {\selectmenu{nombre} à \boolsensing{réponse}}}
\blockvariable{mettre {\selectmenu{nombre} à \ovaloperator{\ovalvariable{nombre} * \ovalnum{1,5}}}}
\blockvariable{ajouter à \selectmenu{nombre} \ovalnum{3}}
\blockvariable{mettre \selectmenu{nombre} à \ovaloperator{\ovalvariable{nombre} * \ovalnum{2}}}
\blockvariable{mettre \selectmenu{nombre} à \ovaloperator{\ovalvariable{nombre} - \ovalnum{6}}}
\blockifelse{si \booloperator{\ovalvariable{nombre} < \ovalnum{20}} alors}
{\blocklook{dire \ovalnum{Bravo, vous avez gagné} pendant \ovalnum{2} seconde}}
{\blocklook{dire \ovalnum{Désolé, vous avez perdu} pendant \ovalnum{2} seconde}}
\end{scratch}
\end{minipage}

\vspace{3mm}
\begin{minipage}{8cm}
\paragraph{ANNEXE}: à coller sur votre copie
\medskip

\begin{scratch}[num blocks]
\blockinit{quand \greenflag est cliqué}
\blockmove{demander \txtbox{bonjour, choisissez un nombre} et attendre}
\blockvariable{mettre {\selectmenu{nombre} à \boolsensing{réponse}}}
\blockvariable{mettre {\selectmenu{nombre} à \ovaloperator{\ovalvariable{nombre} * \ovalnum{1,5}}}}
\blockvariable{ajouter à \selectmenu{nombre} \ovalnum{3}}
\blockvariable{mettre \selectmenu{nombre} à \ovaloperator{\ovalvariable{nombre} * \ovalnum{2}}}
\blockvariable{mettre \selectmenu{nombre} à \ovaloperator{\ovalvariable{nombre} - \ovalnum{6}}}
\blockifelse{si \booloperator{\ovalvariable{nombre} < \ovalnum{20}} alors}
{\blocklook{dire \ovalnum{Bravo, vous avez gagné} pendant \ovalnum{2} seconde}}
{\blocklook{dire \ovalnum{Désolé, vous avez perdu} pendant \ovalnum{2} seconde}}
\end{scratch}
\end{minipage}
\hspace{5mm}
\begin{minipage}{8cm}
\paragraph{ANNEXE}: à coller sur votre copie
\medskip

\begin{scratch}[num blocks]
\blockinit{quand \greenflag est cliqué}
\blockmove{demander \txtbox{bonjour, choisissez un nombre} et attendre}
\blockvariable{mettre {\selectmenu{nombre} à \boolsensing{réponse}}}
\blockvariable{mettre {\selectmenu{nombre} à \ovaloperator{\ovalvariable{nombre} * \ovalnum{1,5}}}}
\blockvariable{ajouter à \selectmenu{nombre} \ovalnum{3}}
\blockvariable{mettre \selectmenu{nombre} à \ovaloperator{\ovalvariable{nombre} * \ovalnum{2}}}
\blockvariable{mettre \selectmenu{nombre} à \ovaloperator{\ovalvariable{nombre} - \ovalnum{6}}}
\blockifelse{si \booloperator{\ovalvariable{nombre} < \ovalnum{20}} alors}
{\blocklook{dire \ovalnum{Bravo, vous avez gagné} pendant \ovalnum{2} seconde}}
{\blocklook{dire \ovalnum{Désolé, vous avez perdu} pendant \ovalnum{2} seconde}}
\end{scratch}
\end{minipage}
\end{document}